\chapter{State of the Art}
\label{chap:state_of_the_art}

The communication and security challenges described in the previous chapter are not specific to renewable energy systems. Similar problems have been studied in the context of industrial IoT and microservice architectures, where researchers have proposed various solutions for multi-protocol communication and API security. Reviewing this literature is essential to understand what design principles and solutions already exist, what their limitations are, and where the contribution of this master's thesis fits within the current state of the art. This chapter organises existing work along four themes: multi-protocol interoperability and translation mechanisms, multi-protocol platform architectures, the Web of Things paradigm as a convergence strategy, and API gateway security in edge and microservice architectures.

% ─────────────────────────────────────────────────────────────────────────────
\section{Protocol Interoperability and Translation}
\label{sec:soa_protocol_translation}
% ─────────────────────────────────────────────────────────────────────────────

Renewable energy devices, like many industrial systems, do not communicate over a single protocol. MQTT, CoAP, HTTP, AMQP, and WebSockets each have legitimate use cases depending on the device type, connectivity constraints, and communication pattern involved. Since protocol convergence on a single standard remains unlikely in the foreseeable future~\cite{derhamy2017}, significant research effort has gone into building translation mechanisms that allow these protocols to coexist within the same system.

Derhamy et al.~\cite{derhamy2017} address this problem in the context of Industrial IoT by proposing a secure, on-demand, and transparent multiprotocol translator built within a Service-Oriented Architecture~(SOA) framework. Their key observation is that direct protocol-to-protocol translation scales quadratically with the number of protocols, following $\frac{n(n-1)}{2}$, which quickly becomes unmanageable. To address this, the authors introduce an intermediate representation \emph{format} (but not an intermediate protocol) that captures all protocol-specific information internally without being transmitted over the wire. This hub-and-spoke architecture reduces the number of required translators to $n-1$ while preserving the information fidelity of direct translation. The system operates on demand, meaning translation instances are provisioned only when a protocol mismatch is detected by the Arrowhead SOA orchestrator, and torn down after a configurable time-to-live. Security is handled through the Arrowhead authorisation system, which issues time-limited tokens and enables fine-grained access control decoupled from the translation logic itself. The translator is implemented in Java and is available as open-source software, actively maintained as part of the broader Eclipse Arrowhead project.

\paragraph{Performance.}
The translator was evaluated on a BeagleBone Black, benchmarked against the Californium HTTP-CoAP proxy running on the same hardware. Round-trip time was measured at the network interface, isolating processing overhead on the device from external network latency. The translator introduced average delays of 50-77~ms depending on payload size, compared to 146-177~ms for the Californium proxy under the same conditions. \\

Da Cruz et al.~\cite{dacruz2019} take a more pragmatic, device-centric approach with \emph{MiddleBridge}, an application-layer gateway that translates MQTT, CoAP, WebSocket, and DDS messages into HTTP. The motivation is that most IoT middleware platforms natively support only HTTP and MQTT, leaving devices using other protocols unable to communicate with the backend without code changes. Unlike earlier bridges such as Ponte and Gothings (both of which have since been deprecated) MiddleBridge additionally applies a payload reduction strategy: rather than forwarding full JSON or XML messages unchanged, it accepts a compact partial payload from the device, reconstructs the full HTTP request internally, and forwards it to the target middleware. The system is deployable on any Java-capable device and exposes a graphical user interface for runtime configuration. MiddleBridge is available as open-source under the GNU~GPL~v3 licence, though its repository has not seen active maintenance since its initial publication.

\paragraph{Performance.}
Experiments were conducted on a Raspberry~Pi~3 hosting MiddleBridge, with an Orion context broker deployed on a separate workstation and load generated via Apache JMeter. Round-trip time and packet size were measured from the IoT device to MiddleBridge only, as MiddleBridge closes the device connection before forwarding to the middleware. Under 100 concurrent users, average round-trip times were 2~ms for MQTT and 8~ms for CoAP, compared to 343~ms for direct HTTP to the broker. In terms of packet size, the payload sent by the IoT device through MiddleBridge was 17 times smaller than an equivalent direct HTTP request.

% ─────────────────────────────────────────────────────────────────────────────
\section{Multi-Protocol IoT Platform Architectures}
\label{sec:soa_multiprotocol_platforms}
% ─────────────────────────────────────────────────────────────────────────────

Beyond point-to-point protocol translation, a parallel line of research has focused on building integrated IoT platforms that natively support multiple application-layer protocols under a unified management and access-control framework.

Akasiadis et al.~\cite{akasiadis2019} present a multi-protocol IoT platform built on top of the SYNAISTHISI cloud platform, integrating four open-source components: RabbitMQ, Ponte, OM2M, and RDF4J, into a single Dockerised deployment that supports MQTT, AMQP, CoAP, REST/HTTP, and WebSockets simultaneously. RabbitMQ and Ponte serve as cross-translation engines, bridging messages across protocols transparently at the messaging layer. The oneM2M common service layer is provided through the OM2M framework, offering standardised device management and resource discovery. Semantic annotations via an RDF4J triplestore enable ontology-based service discoverability, which goes beyond protocol translation and addresses the semantic dimension of interoperability. Authentication and authorisation are implemented through a Flask-based portal, with access control extended to the message brokers so that publish and subscribe permissions are validated dynamically per user and per topic. The platform follows a marketplace architecture promoting service reusability and is delivered as a Dockerised container for local deployment. Of the four constituent components, RabbitMQ and RDF4J remain actively maintained open-source projects; Eclipse Ponte has been archived and is no longer actively developed, and Eclipse OM2M has similarly seen no active maintenance since the early 2020s.

\paragraph{Performance.}
Experiments were run on an Intel Core i5-4440 machine with a separate client machine connected over a 100~Mbps LAN. End-to-end message delivery was measured from publish to reception at a subscribing client. For publish/subscribe protocols, MQTT and AMQP maintained delivery times well below 25~ms under moderate load. AMQP consistently outperformed MQTT, as its authorisation checks are handled natively within RabbitMQ, whereas MQTT relies on external calls to the Flask portal, introducing additional overhead. For request/response protocols (CoAP and HTTP), average translation delays reached into the hundreds of milliseconds due to the polling-based message retrieval model, with CPU usage exceeding 80\% under the same load conditions.

% ─────────────────────────────────────────────────────────────────────────────
\section{The Web of Things Paradigm as an Interoperability Strategy}
\label{sec:soa_wot}
% ─────────────────────────────────────────────────────────────────────────────

The Web of Things~(WoT) paradigm offers a higher-level approach to protocol interoperability by encouraging IoT device resources and functionalities to be exposed as RESTful HTTP services, making them accessible using established web standards regardless of the underlying device protocol.

Benomar et al.~\cite{benomar2024} introduce a gateway-based WoT system that combines two platforms: Stack4Things~(S4T), which handles dynamic DNS and remote accessibility of gateway-hosted services, and the Data eXchange Mediator Synthesizer~(DeXMS), which automatically synthesises protocol mediators that adapt MQTT, CoAP, AMQP, and other protocols into RESTful HTTP interfaces. The system simultaneously addresses two interoperability challenges: protocol heterogeneity, handled by DeXMS mediators, and network reachability, which arises when gateways are deployed behind NATs or firewalls and is addressed by S4T's NGINX reverse proxy and WebSocket tunnel infrastructure. The architecture follows the WoT indirect integration pattern, placing a smart gateway between resource-constrained IoT devices and cloud services, abstracting device protocols and exposing a unified HTTP interface to applications. Mediators are deployed as containerised Java JAR files on the gateway, generated on demand by the cloud-side DeXMS service and managed through IoTronic's remote provisioning system. The system supports bidirectional communication, allowing both data retrieval from devices and actuation towards them. Both Stack4Things and DeXMS are research-oriented open-source platforms whose active maintenance is limited to the academic groups that produced them.

\paragraph{Performance.}
Tests were conducted on a Raspberry~Pi~3 Model~B+ gateway. CPU usage and message throughput were measured per mediator under increasing message rates. For the MQTT-to-HTTP mediator, CPU usage scaled approximately linearly with message rate, reaching around 43\% at 200 messages per second, with message drops beginning near 175 messages per second. For the CoAP-to-HTTP mediator, saturation occurred earlier at approximately 125 messages per second, attributed to the absence of an intermediate message queue. RAM usage remained stable at approximately 9.4\% for MQTT and 9.7\% for CoAP across all tested rates, regardless of message rate.

% ─────────────────────────────────────────────────────────────────────────────
\section{API Gateway Security in Edge and Microservice Architectures}
\label{sec:soa_api_gateway_security}
% ─────────────────────────────────────────────────────────────────────────────

While the works reviewed in the preceding sections primarily address interoperability at the protocol and platform levels, a complementary body of research has focused on the security challenges that arise when IoT devices and microservices are exposed through API gateways in edge computing environments.

Xu et al.~\cite{xu2019} propose a microservice security agent architecture that positions the API gateway as the central security enforcement point in an edge computing deployment. The motivation is a recognised limitation of conventional microservice security models: when security logic is distributed across individual services, each must independently implement and maintain authentication and access control mechanisms, introducing redundancy, inconsistency, and an enlarged attack surface. The proposed architecture addresses this by externalising security concerns to a dedicated agent co-located with the API gateway, so that all inbound requests are intercepted and validated before reaching any backend microservice. The security agent is built on top of the open-source Kong API gateway, deployed within the EdgeX Foundry framework. It integrates two enforcement layers: a JSON Web Token~(JWT) authentication module implementing a Kong plugin that validates HS256- or RS256-signed tokens as specified in RFC~7519, and an Access Control List~(ACL) authorisation module that associates each consumer with a named whitelist group, rejecting requests from consumers not belonging to an authorised group for the target resource. Both layers operate as a unified pipeline, ensuring that JWT authentication is evaluated before ACL authorisation. A companion client server provides a graphical interface for consumer registration, token generation, and ACL group assignment, while a security agent microservice exposes REST APIs to interact with Kong's admin interface programmatically. Kong remains an actively maintained open-source project, and EdgeX Foundry continues to be developed under the Linux Foundation.

\paragraph{Performance.}
Experiments were conducted on a desktop machine running Ubuntu 18.04 with an AMD 64 quad-core processor and 4~GB RAM. Round-trip time was measured from client to microservice for both secured and unsecured paths. Accessing a microservice through the full JWT and ACL enforcement pipeline yielded an average round-trip time of approximately 9~ms, compared to 3~ms for direct access to an unsecured microservice on the same hardware, giving an overhead of approximately 6~ms attributable to the two security enforcement layers.

% ─────────────────────────────────────────────────────────────────────────────
\section{Discussion and Research Gap}
\label{sec:soa_synthesis}
% ─────────────────────────────────────────────────────────────────────────────

The reviewed works collectively address the two main technical dimensions this thesis tackles, yet no single work covers both simultaneously, and none considers the specific operational context of renewable energy asset management.

On the interoperability side, both Derhamy et al.~\cite{derhamy2017} and Da~Cruz et al.~\cite{dacruz2019} demonstrate effective protocol translation mechanisms, but treat security as peripheral to the translation logic. Neither system enforces gateway-level authentication, role-based authorisation, or per-resource access policies. Akasiadis et al.~\cite{akasiadis2019}  and Benomar et al.~\cite{benomar2024} build more complete platform architectures that include basic user authentication, yet neither addresses rate limiting, traffic overload protection, or fine-grained authorization beyond simple credential validation. A further limitation shared by all four systems is that their translation mechanisms operate at the protocol level only: they convert message framing and delivery semantics but do not address device-specific payload formatting. In practice, heterogeneous field devices typically expect command payloads in vendor-specific structures, which pure protocol bridges cannot handle. Our system allows a backend to communicate using a uniformalized format, which then gets translated and transported using device-specific requirements. 

On the security side, Xu et al.~\cite{xu2019}  demonstrate that centralising JWT authentication and ACL-based access control at the gateway layer is both architecturally sound and practically efficient. However, their system operates exclusively over HTTP and assumes well-formed requests from registered consumers, providing no mechanism for handling heterogeneous device protocols or adapting message content.
% I wouldn't put this part:
% Furthermore, none of the reviewed works consider the specific operational constraints of renewable energy assets: proprietary vendor protocols, geographically distributed field deployments, and the requirement for continuous availability in safety-critical installations.

This thesis addresses precisely this gap. The proposed architecture is built on Apache APISIX and integrates, within a single gateway, multi-protocol translation for both HTTP and MQTT devices alongside a comprehensive security model covering JWT-based authentication, role-based authorization, and per-device command access policies. Most importantly, a distinguishing aspect in the translation mechanism of our contribution is that it handles not only protocol conversion but also device-specific payload mapping, allowing backends to issue generic requests that are translated at the gateway into the correct protocol \textbf{and} payload format for each target device. 

The performance characteristics of this architecture are evaluated experimentally in Chapter~\ref{chap:demo_eval}.